% Exam Template for UMTYMP and Math Department courses
%
% Using Philip Hirschhorn's exam.cls: http://www-math.mit.edu/~psh/#ExamCls
%
% run pdflatex on a finished exam at least three times to do the grading table on front page.
%
%%%%%%%%%%%%%%%%%%%%%%%%%%%%%%%%%%%%%%%%%%%%%%%%%%%%%%%%%%%%%%%%%%%%%%%%%%%%%%%%%%%%%%%%%%%%%%

% These lines can probably stay unchanged, although you can remove the last
% two packages if you're not making pictures with tikz.
\documentclass[11pt]{exam}
\RequirePackage{amssymb, amsfonts, amsmath, latexsym, verbatim, xspace, setspace}
\RequirePackage{tikz, pgflibraryplotmarks}

% By default LaTeX uses large margins.  This doesn't work well on exams; problems
% end up in the "middle" of the page, reducing the amount of space for students
% to work on them.
\usepackage[margin=1in]{geometry}
\usepackage{enumerate}
\usepackage{amsthm}

\theoremstyle{definition}
\newtheorem{soln}{Solution}

% Here's where you edit the Class, Exam, Date, etc.
\newcommand{\class}{Math 150A Section 8 and 16}
\newcommand{\term}{Spring 2025}
\newcommand{\examnum}{Practice Exam II Solutions}
\newcommand{\examdate}{March 06, 2025}
\newcommand{\timelimit}{1 Hour 50 Minutes}

% For an exam, single spacing is most appropriate
\singlespacing
% \onehalfspacing
% \doublespacing

% For an exam, we generally want to turn off paragraph indentation
\parindent 0ex

\begin{document} 

% These commands set up the running header on the top of the exam pages
\pagestyle{head}
\firstpageheader{}{}{}
\runningheader{\class}{\examnum\ - Page \thepage\ of \numpages}{\examdate}
\runningheadrule

\begin{flushright}
\begin{tabular}{p{2.8in} r l}
\textbf{\class} & \textbf{Name (Print):} & \makebox[2in]{\hrulefill}\\
\textbf{\term} &&\\
\textbf{\examnum} & \textbf{Student ID:}&\makebox[2in]{\hrulefill}\\
\textbf{\examdate} &&\\
\textbf{Time Limit: \timelimit} % & Teaching Assistant & \makebox[2in]{\hrulefill}
\end{tabular}\\
\end{flushright}
\rule[1ex]{\textwidth}{.1pt}


This exam contains \numpages\ pages (including this cover page) and
\numquestions\ problems.  Check to see if any pages are missing.  Enter
all requested information on the top of this page, and put your initials
on the top of every page, in case the pages become separated.\\

You may \textit{not} use your books or notes on this exam.  However, you may use a single, handwritten, one-sided notesheet and a \textit{basic} calculator.\\

You are required to show your work on each problem on this exam.  The following rules apply:\\

\begin{minipage}[t]{3.7in}
\vspace{0pt}
\begin{itemize}

%\item \textbf{If you use a ``fundamental theorem'' you must indicate this} and explain
%why the theorem may be applied.

\item \textbf{Organize your work}, in a reasonably neat and coherent way, in
the space provided. Work scattered all over the page without a clear ordering will 
receive very little credit.  

\item \textbf{Mysterious or unsupported answers will not receive full
credit}.  A correct answer, unsupported by calculations, explanation,
or algebraic work will receive no credit; an incorrect answer supported
by substantially correct calculations and explanations might still receive
partial credit.  This especially applies to limit calculations.

\item If you need more space, use the back of the pages; clearly indicate when you have done this.

\item \textbf{Box Your Answer} where appropriate, in order to clearly indicate what you consider the answer to the question to be.
\end{itemize}

Do not write in the table to the right.
\end{minipage}
\hfill
\begin{minipage}[t]{2.3in}
\vspace{0pt}
%\cellwidth{3em}
\gradetablestretch{2}
\vqword{Problem}
\addpoints % required here by exam.cls, even though questions haven't started yet.	
\gradetable[v]%[pages]  % Use [pages] to have grading table by page instead of question

\end{minipage}
\newpage % End of cover page

%%%%%%%%%%%%%%%%%%%%%%%%%%%%%%%%%%%%%%%%%%%%%%%%%%%%%%%%%%%%%%%%%%%%%%%%%%%%%%%%%%%%%
%
% See http://www-math.mit.edu/~psh/#ExamCls for full documentation, but the questions
% below give an idea of how to write questions [with parts] and have the points
% tracked automatically on the cover page.
%
%
%%%%%%%%%%%%%%%%%%%%%%%%%%%%%%%%%%%%%%%%%%%%%%%%%%%%%%%%%%%%%%%%%%%%%%%%%%%%%%%%%%%%%

\begin{questions}

\addpoints

\question[10]\mbox{}
\textbf{TRUE or FALSE!}  Write 
\begin{enumerate}[(a)]
\item if $f''(3) = 0$ then $f(x)$ has an inflection point at $3$
\vspace{1.2in}
\item the function $f(x) = |x|$ has no critical points
\vspace{1.2in}
\item if $f(x)$ is continuous on the interval $[0,1]$, then it has a has both absolute maximum and absolute minimum values on $[0,1]$
\vspace{1.2in}
\item if $f'(x) < 0$ for all $x$, then $f(x)$ can have at most one root
\vspace{1.2in}
\item if $f(0)=0$ and $f(1)=60$ then $f'(x)=60$ for some value of $x$ in $(0,1)$
\vspace{1.2in}
\end{enumerate}

F,F,T,T,T

\newpage
\question[10]\mbox{}
Water is leaking out of an inverted conical tank at a rate of $10,000$ cm$^3$/min at the same time water is being pumped into the tank at a constant rate.
The tank has a height of $6$ m and the diameter at the top is $4$ m.
If the water level is rising at a rate of $20$ cm/min when the height of the water is $2$ m, what is the rate at which water is being pumped into the tank?

\textbf{Solution:} 
Always draw a figure first!  However, making one on the computer here is inconvenient, so we will skip this step.\\
Let $V$ be the volume of the tank, $h$ be the height of the water level, and $d$ be the diameter of the water's surface.
Then we know $\frac{dh}{dt} = 20$ cm/min and we want the rate that water goes into the tank.
We will get this by calculating the rate of change of the volume in the tank $\frac{dV}{dt}$.

The volume of water in the tank is related to the diameter and the height by
$$V = \frac{1}{12}\pi d^2h.$$
Also, by similar triangles we know that
$$\frac{d}{h} = \frac{4}{6}.$$
Therefore $d = 2h/3$, and
$$V = \frac{1}{9}\pi h^3.$$
Thus by related rates
$$\frac{dV}{dt} = \frac{dV}{dh}\frac{dh}{dt} = \frac{1}{3}\pi h^2\cdot 20.$$
At the point when $h=2$ meters, ie. $200$ cm, this says $\frac{dV}{dt} = \frac{800000\pi}{3}$.
Now, the change in the volume of the tank is the rate at which water goes in minus the rate at which water goes out, so
$$\frac{800000\pi}{3} = \frac{dV}{dt} = \text{(Rate in)} - \text{(Rate out)} = \text{(Rate in)} - 10000.$$
Therefore
$$\text{(Rate in)} = \frac{800000\pi}{3} + 10000\ \text{cm$^3$/min}.$$

\newpage
\question[6]
In class, we learned a list of critical properties to consider when sketching a function.
Suppose that $f(x)$ is a function which simultaneously satisfies all of the following:
\begin{enumerate}[(A)]
\item the domain is $\mathbb{R}$
\item there is a $y$-intercept at $0$; there are $x$-intercepts at $0$ and $\pm 3\sqrt{3}$
\item the function is odd
\item there are no horizontal or vertical asymptotes
\item the function is increasing on $(-\infty,-1)$ and $(1,\infty)$, and is decreasing on $(-1,1)$
\item it has a local maximum $f(-1)=2$ and local minimum $f(1)=-2$
\item the function is concave upward on $(0,\infty)$, concave downward on $(-\infty,0)$ with an inflection point at $x=0$
\end{enumerate}
Using these properties, we should be able to produce a remarkably good sketch of $f(x)$, even though we don't even know what $f(x)$ is!  Create a sketch of $f(x)$ which simultaneously satisfies all the above properties.

\textbf{Solution:} OMITTED

\newpage
\question[10]\mbox{}
\begin{enumerate}[(a)]
\item  State the Mean Value Theorem
\vspace{2in}
\item  State the Extreme Value Theorem
\end{enumerate}
\textbf{Solution:}  
\begin{enumerate}[(a)]
\item  MVT says that if $f(x)$ is continuous on $[a,b]$ and differentiable on $(a,b)$, then there exists $c\in (a,b)$ with
$$f'(c) = \frac{f(b)-f(a)}{b-a}$$
\vspace{2in}
\item  The Extreme Value Theorem says that if $f(x)$ is continuous on $[a,b]$ then the absolute maximum and absolute minimum values of $f(x)$ on $[a,b]$ exist.
\end{enumerate}



\newpage
\question[10]\mbox{}
For each of the following, provide an example or explain why it does not exist.
\begin{enumerate}[(a)]
\item  a function $f(x)$ with a critical point at $0$ but neither a local max nor local min at $x=0$
\vspace{1.2in}
\item  a function $f(x)$ which satisfies $f''(1) = 0$ but which does not have an inflection point at $x=1$
\vspace{1.2in}
\item  a nonconstant function $f(x)$ whose derivative is always zero
\vspace{1.2in}
\item  a function with an inflection point at $x=0$
\vspace{1.2in}
\item  a differentiable function with three $x$-intercepts, whose derivative only has one $x$-intercept
\vspace{1.2in}
\end{enumerate}
\textbf{Solution:}
\begin{enumerate}[(a)]
\item  $f(x) = x^3$
\item  $f(x) = (x-1)^4$
\item  can't exist because of MVT
\item  $f(x) = x^3$
\item  can't exist because of MVT
\end{enumerate}

\newpage
\question[10]
If $1,200$ cm$^2$ material is available to make a box with a square base and an open top, find the largest possible volume of the box.

\textbf{Solution:}
Let $x$ be the width of the box and $y$ be the height.
The surface area of the box is $4xy+x^2$, so if we use up all our material we must have
$$4xy+x^2=1200.$$
This means
$$y = \frac{300}{x} - \frac{1}{4}x.$$
The volume of the box is
$$V = x^2y,$$
so subbing in the previous equation we get
$$V = 300x - \frac{1}{4}x^3.$$
The critical points of this are $x=\pm 20$, and the absolute max occurs at $x=20$.
At this value, $V = 4000$ so the largest possible volume is $4000$ cubic centimeters.

\newpage
\question[7]
Consider the function
$$f(x) = x^4 +x^2-1$$
\begin{enumerate}[(a)]
\item show that $f(x)$ has only one critical point
\vspace{2in}
\item explain why the Mean Value Theorem says $f(x)$ cannot have more than two roots
\vspace{2in}
\item calculate the value of $f(-1)$, $f(0)$ and $f(1)$
\vspace{2in}
\item explain why the Intermediate Value Theorem says $f(x)$ has at least two real roots (and therefore precisely two roots)
\end{enumerate}
\textbf{Solution:}
\begin{enumerate}[(a)]
\item $f'(x) = 4x^3+2x = 2x(2x^2+1)$.  The only place where this is zero is $x=0$, so there is only one critical point.
\item If $f(x)$ has three roots at $x=a$, $x=b$, and $x=c$ (in increasing order), then the Mean Value Theorem (or Rolle's Theorem) says that $f'(x)=0$ at some point between $a$ an $b$.  Likewise, $f'(x) = 0$ at some point between $b$ and $c$.  However, there's only one point where $f'(x)=0$, so that's not possible!
\item $f(-1) = 1$, $f(0) = -1$, and $f(1) = 1$
\item IVT says that $f$ has a zero between $-1$ and $0$.  Likewise, MVT says $f$ has a zero between $0$ and $1$.
\end{enumerate}


\end{questions}
\end{document}
